\chapter{Acceleration Searching} \label{appen:zparams}

Searching for binary pulsars requires a slightly different approach as the motion of the system causes the observed pulse frequency to smear across the Fourier bins \citep{ng_high_2015}, in turn this reduces the sensitivity of the search. A solution to this is to split up the search into smaller time intervals and assume that the radial velocity is a constant on this timescale. This can be shown to be a good approximation for $P_b/10$. \\ 
The spin frequency, $f_{\text{spin}}$ and the time of pulse emission, $t_{\text{pulse}}$ the pulsar's phase can be expressed as, $\phi_p = f_{\text{spin}}t_{\text{pulse}}$. Moreover, the pulse time can be expressed as the time of arrival from the pulsar, $d(t_0)$, 

\begin{equation}
    t_{\text{pulse}} = t_0 - \frac{d(t_0)}{c}
\end{equation}

thus the phase can be expressed as,

\begin{align}
    \phi_p = & f_{\text{spin}}\left[ t_0 - \frac{d(t_0)}{c}\right] \\ 
    = & f_\text{spin} \left[ t_0 - \frac{a \sin i}{c} \sin \left[ \frac{2 \pi (t - t_{\text{asc}})}{P_B} \right] \right]
\end{align}

where $a$ is the semi-major axis, $i$ is the inclination angle, $t_{\text{asc}}$ is the time of ascending node and $P_B$ is the orbital period. A Taylor expansion can be used to $\sin x$ around $x = a$ such that, 

\begin{equation}
    \sin (x) \simeq   \sin (a) + \cos (a)(x - a) - \frac{\sin (a)}{2!}(x - a)^2 - \frac{\cos (a)}{3!}(x - a)^3 + \ldots
\end{equation}

Therefore the phase of the pulsar with a constant spin down rate can be expressed as, 

\begin{equation}
  \phi(t) = f^1 t_0 + \frac{\dot f}{2} (t - t_0)^2 + \frac{\ddot f}{6} (t - t_0)^3 + \ldots 
\end{equation}

Matching coefficients between the full Taylor expansion and the phase expression for the pulsar gives the following relation, $\dot f/2 = f_{\text{spin}} \frac{a \sin i}{c} A$. Where $A$ is represents all the terms in the expansion. The final simplified expression works out to be, 

\begin{equation}
    \dot f = f_\text{spin} \frac{4 \pi^2 a \sin i }{cP^2_B} \sin \left[ \frac{2 \pi (t - t_{\text{asc}})}{P_B} \right]
\end{equation}

Thus over a small enough time period spin-down is approximately constant.

\begin{equation}
    \frac{k_2 P_B}{2 \pi q} = a \sin i   
    \label{eq:constant_accel}
\end{equation}

Where $k_2$ is the radial velocity semi-amplitude, $P_B$ is the orbital period and $q$ is the mass ratio.

Searching is carried out using \texttt{accelsearch} which searches the FFT time series for period emissions by summing the highest Fourier frequency derivative. The number of bins that this signal is smeared is given by the acceleration parameter \citep{ransom_new_2001}, 

\begin{equation}
    z \simeq \dot f T^2_{\text{obs}}
\end{equation}

Where $T_{\text{obs}}$ is the length of the observation. The search is carried out over a range of accelerations which can be based on \cref{eq:constant_accel} if the physical quantities are known. This is the means with informed the $z$-parameters used throughout \cref{chap:redback}